%!TeX encoding = utf8
%!TeX spellcheck = it_IT

\documentclass[a4paper, 11pt]{book}

\usepackage[T1]{fontenc}
\usepackage[utf8]{inputenc}
\usepackage[italian]{babel}
\usepackage{amsmath, amssymb, amsthm}
\usepackage{geometry, emptypage}
\usepackage{graphicx, subfig}
\usepackage{enumitem, lipsum}
\usepackage{marginnote}
%\usepackage[autostyle,italian=guillemets]{csquotes}
%\usepackage[style=philosophy-modern, backend=biber, square, hyperref, backref]{biblatex}
\usepackage{listings}
\usepackage[eulerchapternumbers, palatino=false, parts=false]{classicthesis}
\usepackage{arsclassica}

% Font
%\IfFileExists{MinionPro.sty}{
%	\usepackage[noopticals, mathlf, minionint, loosequotes]{MinionPro}
%	\usepackage[toc, enum]{tabfigures}
%	\usepackage{relsize}
%	\let\oldsum\sum\let\oldprod\prod
%	\def\sum{\mathlarger{\oldsum}}%
%	\def\prod{\mathlarger{\oldprod}}
%}{
	\usepackage[proportional, oldstyle]{cochineal}
	\usepackage[cochineal, vvarbb]{newtxmath}
%}

% Figures
\graphicspath{{Figure/}}
\captionsetup{format=hang}

% Layout
\geometry{%showframe,
	width=30pc, height=54pc, vmarginratio=1000:1414}
\setlength\parindent{1em}
\newcommand\setupheading[1]{%
	\cleardoublepage\phantomsection
	\markboth{\spacedlowsmallcaps{#1}}{\spacedlowsmallcaps{#1}}%
}

% Bibliography
%\addbibresource{Bibliography.bib}
%\renewcommand*{\nameyeardelim}{\addcomma\space}

% Headings
\renewcommand{\sectionmark}[1]{
	\markright{\thesection\enspace\textit{#1}}
} 
\clearplainofpairofpagestyles

% Code
\lstset{language=C,
 	basicstyle=\ttfamily\small,
	keywordstyle=\color{RoyalBlue},
	commentstyle=\color{black!20!Emerald}\footnotesize\itshape,
	columns=flexible,
	aboveskip=\bigskipamount,
	stringstyle=\color{Orange},
	numbers=none,
	tabsize=4,
	showstringspaces=false,
	breaklines=true,
	frameround=ftff,
	frame=lines,
	backgroundcolor=\color{white},%\color[gray]{.99},
	escapeinside={£}{£},
	extendedchars=true,
	literate={à}{{\`a}}1 {è}{{\`e}}1 {ù}{{\`u}}1 {ö}{{\"o}}1
}

% Theorems
\newtheoremstyle{plain}{}{}{\itshape}{}{}{}{ }{%
	\textbf{\thmname{#1}\thmnumber{ #2.}}\thmnote{ (#3).}}
\newtheoremstyle{definition}{}{}{}{}{}{}{ }{%
	\textbf{\thmname{#1}\thmnumber{ #2.}}\thmnote{ (#3).}}

\theoremstyle{plain}
\newtheorem{theorem}{Theorema}[chapter]
\newtheorem{corollary}[theorem]{Corollario}
\newtheorem{proposition}[theorem]{Proposizione}
\newtheorem{lemma}[theorem]{Lemma}

\theoremstyle{definition}
\newtheorem{definition}[theorem]{Definizione}
\newtheorem{remark}[theorem]{Osservazione}
\newtheorem{example}[theorem]{Esempio}

\numberwithin{equation}{chapter}
\numberwithin{figure}{chapter}
\numberwithin{table}{chapter}

% Other
\hypersetup{%hidelinks,
	linkcolor=RoyalBlue}
\microtypesetup{kerning=true}
\let\marginpar\marginnote
\setlist{itemsep=0pt}

% Dates
\usepackage{Files/datemath}
\newdate{lecture}
\setdaynames{domenica}{luned\`i}{marted\`i}{mercoled\`i}{gioved\`i}{venerd\`i}{sabato}
\setmonthnames{gennaio}{febbraio}{marzo}{aprile}{maggio}{giugno}{luglio}
	{agosto}{settembre}{ottobre}{novembre}{dicembre}
\newcommand\lecture[1]{%
	\setdate{lecture}{#1/2019}%
	\calcdayofweek{lecture}%
	\section{Lezione di \dayofweek{lecture} \theday{lecture} \monthname{lecture}}%
}

\begin{document}

\mainmatter

% !TeX source = ../Tesi.tex

%****************************************************************
% Title Page
%****************************************************************

\begin{titlepage}
\newgeometry{top=1in, bottom=1in, left=1.5in, right=1.5in}
\begin{center}
\includegraphics[width=3cm]{LogoUnipi}\par\bigskip
%\includegraphics[width=5cm]{ScrittaUnipi}\par\bigskip
{\large\spacedallcaps{Universit\`a di Pisa}}\\[1ex]
{\large Dipartimento di Informatica}

\vspace{\stretch{1}}

{\Large\spacedallcaps{Machine Learning}}\par\bigskip
{\large\itshape Appunti dalle lezioni del Prof. Alessio Micheli}

\vspace{\stretch{1}}

\begin{center}
\itshape
\begin{tabular}{r@{\hspace{.5em}}l}
Pasquale & Miglionico\\[1.5ex]
Guido & Narduzzi\\[1.5ex]
Enrico & Negri\\[1.5ex]
Filippo & Quattrocchi\\[1.5ex]
Francesco & Zigliotto
\end{tabular}
\end{center}

\vspace{\stretch{1}}

\spacedlowsmallcaps{Anno accademico 2019--2020}
\end{center}
\restoregeometry
\end{titlepage}

% !TeX source = ../Tesi.tex

%****************************************************************
% Contents
%****************************************************************

\begingroup
\manualmark
\setupheading{\contentsname}
\tableofcontents
\endgroup

\chapter{Esempi}

\lipsum[2]

%!TeX encoding = utf8
%!TeX spellcheck = it_IT
%!TeX source = ../ML.tex

\lecture{22/9}

\subsection{Alcune indicazioni}

Un paio di cose:
\begin{itemize}
\item Per creare paragrafi numerati, non va usato \lstinline|\section|, ma \lstinline|\subsection|;
\item Non usate il grassetto, ma il \emph{corsivo};
\item Usate \texttt{equation}  per le equazioni, in modo che vengano tutte numerate;
\item usare gli ambienti \lstinline|definition|, \lstinline|theorem|, \lstinline|example|.
\end{itemize}

\begin{definition}[Tensore]
Un \emph{tensore} è ciò che ruota come un tensore.
\end{definition}

\begin{equation}
e^{z}=\sum_{n=0}^{+\infty} \frac{z^n}{n!}.
\end{equation}

\lipsum[3]

\begin{lstlisting}[language=C]
#include <stdio.h>
int main()
{
   // printf() displays the string inside quotation
   printf("Hello, World!");
   return 0;
}
\end{lstlisting}

\end{document}